\documentclass[10pt,a4paper]{amsart}
\usepackage[margin=19mm]{geometry}
\usepackage{amsmath,amssymb,mathtools}
\usepackage[scr=rsfs,cal=euler]{mathalfa}
\usepackage{xfrac}
\usepackage[unicode,hidelinks,bookmarks=false]{hyperref}
\newcommand{\ie}{i.e.\ }
\newcommand{\cf}{cf.\ }
\newcommand{\resp}{respectively\ }
\newcommand{\Subsection}{\subsection}
\providecommand{\nobreakdash}{\nobreak\hbox{-}}
\setlength{\emergencystretch}{2em}
\multlinegap0pt
\usepackage{mathtools,mathabx}
\usepackage{enumitem}
\usepackage{cleveref}
\usepackage{aliascnt}
\newtheorem{theorem}{Theorem}[section]
\newaliascnt{lemma}{theorem}
\newtheorem{lemma}[lemma]{Lemma}
\newaliascnt{proposition}{theorem}
\newtheorem{proposition}[proposition]{Proposition}
\newaliascnt{corollary}{theorem}
\newtheorem{corollary}[corollary]{Corollary}
\newaliascnt{definition}{theorem}
\newtheorem{definition}[definition]{Definition}
\newaliascnt{example}{theorem}
\newtheorem{example}[example]{Example}
\newaliascnt{assumption}{theorem}
\newtheorem{assumption}[assumption]{Assumption}
\newaliascnt{remark}{theorem}
\newtheorem{remark}[remark]{Remark}
\allowdisplaybreaks[4]
\DeclareMathOperator*{\sgn}{sgn}
\DeclareMathOperator*{\Cov}{Cov}
\DeclareMathOperator*{\Var}{Var}
\DeclareMathOperator*{\argmin}{argmin}
\providecommand{\supp}{\operatorname{supp}}
\newcommand{\Pp}{\mathbb P}
\newcommand{\Ee}{\mathbb E}
\newcommand{\cH}{\mathcal H}
\newcommand{\cM}{\mathcal M_1}
\begin{document}
\title[High Minima of Gaussian Processes: Overshoots and Minimizer ]{High Minima of Gaussian Processes: Overshoots and Minimizer Locations}
\author{Secondary 60G70, 60G17, 60G22; Enkelejd Hashorva; Svyatoslav Novikov}
\date{}
\maketitle
\noindent\emph{Source and licence.} High Minima of Gaussian Processes: Overshoots and Minimizer Locations. Version arXiv:2607.20714v1 (CC BY 4.0). This material is distributed under the Creative Commons Attribution 4.0 International licence, http://creativecommons.org/licenses/by/4.0/, as recorded in the arXiv OAI licence metadata for this exact version and shown on the abstract page https://arxiv.org/abs/2607.20714v1. Full rights notice: licenses/arxiv-2607.20714-CC-BY-4.0.txt.\par\smallskip
\noindent\textit{Editorial scope.} Counted unit: the original \texttt{theorem} environment \texttt{thm:optimal-measure-criterion} at source lines 591--617 together with its complete proof at lines 619--652; the delivered document keeps source lines 550--653 so that every referenced label and every citation resolves inside it. Native editable source is the paper's own concatenated TeX; the AMS wrapper is editorial formatting only. No publisher class, upstream script or unrelated artwork is redistributed.
\section{Proofs}
\label{sec:proofs}
We first present the RKHS identities for the covariance energy and
potential defined above.  Let $\cH$ be the RKHS of the covariance kernel
$R$, with inner product $\langle\cdot,\cdot\rangle_{\cH}$ and norm
$\|\cdot\|_{\cH}$.  Since
\[
 \|R(s,\cdot)-R(t,\cdot)\|_{\cH}^2
 =R(s,s)+R(t,t)-2R(s,t), \qquad s,t \in K
\]
the map $t\mapsto R(t,\cdot)$ is continuous from $K$ into $\cH$.
Consequently, for every finite signed measure $\eta$ on $K$, the
Bochner integral
\[
 k_\eta=\int_KR(t,\cdot)\,\eta(dt)
\]
belongs to $\cH$.  Define the isonormal map first on kernel sums by
\[
 W\left(\sum_{j=1}^nc_jR(t_j,\cdot)\right)
 =\sum_{j=1}^nc_jX(t_j)
\]
and extend it by isometry to $\cH$.  Approximation by kernel sums and
Fubini's theorem, justified by Fernique integrability, then give
\begin{equation}
 \begin{split}
 W(k_\eta)&=\int_KX(t)\,\eta(dt),\\
 \langle h,k_\eta\rangle_{\cH}&=\int_Kh(t)\,\eta(dt),\qquad h\in\cH,\\
 \|k_\eta\|_{\cH}^2
 &=\mathcal E_R(\eta)
 =\Ee\left\{\left(\int_KX(t)\,\eta(dt)\right)^2\right\}\geq0.
 \end{split}
 \label{eq:rkhs-measure-identities}
\end{equation}
In particular, every $h\in\cH$ is continuous, hence bounded, on $K$,
because $h(t)=\langle h,R(t,\cdot)\rangle_{\cH}$.

We shall use the following standard optimality criterion for minimum-energy
probability measures; see   \cite[Theorem~4.3]{adler2014existence}, \cite[Theorem~2.1]{selk2024large} and
\cite[Theorem~1.1]{pronzato2021minimum}.  We include its short proof for
completeness.


\begin{theorem}
\label{thm:optimal-measure-criterion} The probability measure
\(\mu\in\cM(K)\) is optimal
if and only if

\[
        k_\mu(t)\ge \mathcal E_R(\mu),
        \qquad t\in K.
\]

Moreover, if \(\mu\) is optimal, then
\[
        k_\mu(t)=\mathcal E_R(\mu),
        \qquad t\in\supp(\mu).
\]

In particular, if \(k_\mu\) is constant on \(K\), then \(\mu\) is
optimal. If, in addition,
\[
        \mathcal E_R(\eta)>0
\]
for every nonzero finite signed measure \(\eta\) on \(K\) satisfying
\(\eta(K)=0\), then the optimal measure is unique.

Consequently, if \(\supp(\mu)=K\), then \(\mu\) is optimal if and only if
\(k_\mu\) is constant on \(K\).
\end{theorem}

\begin{proof}
For \(\nu\in\cM(K)\), put \(\eta=\nu-\mu\).  For
\(0\leq\varepsilon\leq1\),
\begin{align*}
 \mathcal E_R(\mu+\varepsilon\eta)
 &=\mathcal E_R(\mu)
   +2\varepsilon\left(\int_Kk_\mu(t)\,\nu(dt)
                         -\mathcal E_R(\mu)\right)
   +\varepsilon^2\mathcal E_R(\eta).
\end{align*}
If \(\mu\) is optimal, the right derivative at zero is nonnegative.
Taking \(\nu=\delta_t\) gives
\(k_\mu(t)\geq\mathcal E_R(\mu)\) for every \(t\in K\).  Since
\[
 \int_Kk_\mu(t)\,\mu(dt)=\mathcal E_R(\mu),
\]
equality holds \(\mu\)-almost everywhere and, by continuity, on
\(\supp(\mu)\).

Conversely, suppose that
\(k_\mu\geq\mathcal E_R(\mu)\) on \(K\).  The same expansion with
\(\varepsilon=1\) gives
\[
 \mathcal E_R(\nu)-\mathcal E_R(\mu)
 =2\left(\int_Kk_\mu(t)\,\nu(dt)-\mathcal E_R(\mu)\right)
   +\mathcal E_R(\nu-\mu)\geq0,
\]
and hence \(\mu\) is optimal.  If \(k_\mu\) is constant, its constant
value is \(\mathcal E_R(\mu)\), so the criterion applies.  Strict positivity
of \(\mathcal E_R\) on nonzero signed measures of total mass zero makes
\(\mathcal E_R\) strictly convex on \(\cM(K)\), and therefore the
optimal measure is unique.  The last assertion follows from the equality
on \(\supp(\mu)\).
\end{proof}


\begin{thebibliography}{99}
\bibitem[adler2014existence]{adler2014existence} R.~J. Adler, E.~Moldavskaya, and G.~Samorodnitsky, ``On the existence of paths between points in high level excursion sets of {G}aussian random fields,'' {\em The Annals of Probability}, vol.~42, no.~3, pp.~1020--1053, 2014.
\bibitem[pronzato2021minimum]{pronzato2021minimum} L.~Pronzato and A.~Zhigljavsky, ``Minimum-energy measures for singular kernels,'' {\em Journal of Computational and Applied Mathematics}, vol.~382, p.~113089, 2021.
\bibitem[selk2024large]{selk2024large} Z.~Selk, ``Large deviations for high minima of {G}aussian processes with nonnegatively correlated increments,'' {\em Statistics \& Probability Letters}, vol.~206, p.~110001, 2024.
\end{thebibliography}
\end{document}
